\documentclass{article}

\PassOptionsToPackage{numbers, compress}{natbib}

\usepackage[preprint]{neurips}

\usepackage[utf8]{inputenc} % allow utf-8 input
\usepackage[T1]{fontenc}    % use 8-bit T1 fonts
\usepackage{hyperref}       % hyperlinks
\usepackage{url}            % simple URL typesetting
\usepackage{booktabs}       % professional-quality tables
\usepackage{amsfonts}       % blackboard math symbols
\usepackage{nicefrac}       % compact symbols for 1/2, etc.
\usepackage{microtype}      % microtypography
\usepackage{xcolor}         % colors
\usepackage{amsmath}
\usepackage{graphicx}
\usepackage{enumitem}
\usepackage{xspace}

\title{Circuits of Uncertainty in Language Models}

\author{%
  Karan Singh \\
  % Dept. of Electrical Engineering \\
  Stanford University \\
  \And
  Chelsea Zou \\
  % Symbolic Systems Program \\
  Stanford University
}


\begin{document}
\maketitle

\begin{abstract}
Language models often answer incorrectly despite expressing high confidence. We investigate whether hidden-state representations contain a correctness signal that is separable from confidence, where this signal emerges, and whether it causally contributes to answer generation. Across five QA datasets and multiple Qwen and Llama models, linear probes decode correctness from residual-stream activations with AUCs above 0.80, substantially outperforming verbalized confidence. The correctness and confidence directions are nearly orthogonal in activation space. Activation patching shows that late-layer representations, especially MLP outputs, causally increase next-token evidence for correct answers, and this effect replicates across models. Yet these interventions rarely alter final predictions, and steering along a 1D correctness direction fails to improve correctness or confidence. Our findings reveal a distinction between decodable correctness signals, logit-causal representations, and representations that can be leveraged for behavioral control.
\end{abstract}

\section{Introduction}

Language models can confidently produce factual errors. Yet a growing body of work suggests that this apparent overconfidence is incomplete: hidden-state representations often contain information about whether an answer is actually correct, even when the model’s verbal confidence suggests otherwise \citep{Azaria2023TheLying,Burns2024DiscoveringSupervision,Marks2024TheDatasets}. In other words, a model may internally encode evidence that an answer is unreliable while still producing that answer. This raises a fundamental question: what is the relationship between what language models say and what they know? If correctness is represented internally, where does that representation emerge, and does it play a causal role in answer generation?

Answering these questions requires distinguishing three increasingly strong notions of internal knowledge. A correctness signal may be decodable, meaning it can be recovered from hidden activations. It may be causally relevant, meaning interventions on that representation alter next-token predictions. Or it may be behaviorally controllable, meaning manipulating the signal reliably changes the final answer. These notions need not coincide: a representation can be highly predictive without participating in computation, and it can influence logits without being sufficient to change behavior.

In this work, we systematically investigate the relationship between correctness, confidence, and causality in modern language models. Our contributions can be summarized as follows:

\begin{itemize}[topsep=0pt,leftmargin=4mm,itemsep=0pt,partopsep=0pt,parsep=2pt]
    \item We construct a probing benchmark for factual question answering across five datasets and multiple Qwen and Llama models. We evaluate correctness probes throughout the residual stream and compare them against log-probability, entropy, and verbal-confidence baselines.
    % \item We show that correctness is substantially more decodable than confidence-based signals would suggest, while verbalized confidence is consistently weaker and nearly orthogonal to the correctness direction in activation space.
    \item We localize the representations most responsible for this signal using activation patching at the residual-stream, component, head, position, and question levels. 
    \item We demonstrate a gap between prediction, causation, and control. Although correctness is decodable and partially causal at the logit level, interventions along a one-dimensional correctness direction rarely change final answers and do not improve correctness or confidence.
\end{itemize}

Taken together, our results reveal three distinct notions of internal knowledge in language models: information that can be read out from hidden states, information that causally influences token predictions, and information that can be leveraged to reliably control model behavior. These notions are often treated as interchangeable, but our findings show that they can diverge sharply in practice.

\begin{figure}[t]
    \centering
    \includegraphics[width=\textwidth]{figures/intro.png}
    \caption{\textit{Where do language models internally know that they are wrong?} We evaluate this question using open-ended factual QA across five datasets. For each generation, we extract residual-stream activations throughout the network and train linear probes to predict whether the final answer is correct. We compare these internal signals against output-level confidence measures such as entropy, verbalized confidence, self-consistency, and answer log-probability. The resulting benchmark isolates hidden-state representations of correctness from confidence expressed in model outputs.}
    \label{fig:intro}
\end{figure}

\section{Related Work}

\paragraph{Model confidence and calibration.}

Calibration work studies whether predicted confidence tracks empirical correctness, from classic post-hoc calibration in modern neural networks to language-model-specific analyses of generative question answering and few-shot prompting \cite{Guo2017Calibration,jiang-etal-2021-know,zhao2021calibrate}. In large language models, useful output-level signals include token log-probabilities, entropy and margins, self-evaluation scores, verbalized confidence, and sample- or semantics-based uncertainty measures such as self-consistency, SelfCheckGPT, semantic uncertainty, and semantic entropy \cite{Kadavath2022LanguageKnow,Lin2022TruthfulQA:Falsehoods,manakul2023selfcheckgpt,tian2023just,xiong2024can}. These measures are attractive because they require only model outputs, but they remain incomplete: models can still imitate fluent falsehoods, as emphasized by TruthfulQA, and verbalized confidence captures what a model says about its uncertainty rather than necessarily the internal representations that drive generation \cite{Lin2022TruthfulQA:Falsehoods}. We therefore treat output-level uncertainty and confidence as strong baselines and ask when hidden states add information beyond them. 

\paragraph{Probing hidden representations.}

Linear probes have long been used to test what information is encoded in neural activations, from early linear classifier probes to structural and edge probes for linguistic structure \cite{Alain2017UnderstandingRepresentations, Hewitt2019DesigningProbes, tenney2019you}. At the same time, the probing literature has repeatedly warned that decodability alone does not establish mechanism: control tasks, selectivity analyses, information-theoretic formulations, and explicit counterexamples show that a probe may exploit correlations or information not actually used by the model for its task \cite{Hewitt2019DesigningProbes, pimentel-etal-2020-information, ravichander-etal-2021-probing, belinkov2022probing}. Against this backdrop, recent work has shown that truthfulness and non-factuality are often detectable in hidden states, including unsupervised latent-knowledge probes, truthfulness classifiers for generated statements, linear “truth directions,” settings where probes recover knowledge despite systematically wrong outputs, and detectors based on activation geometry or question representations \cite{Burns2024DiscoveringSupervision,Azaria2023TheLying, Marks2024TheDatasets,mallen2023not}. We build on this work by focusing on short-form factual QA and testing whether the same signal is also causal. 

\paragraph{Mechanistic interpretability and causal interventions.}
Mechanistic interpretability seeks to localize the computations underlying model behavior, from the circuits program and transformer-circuit analyses to causal mediation, interchange interventions, and causal-abstraction frameworks \cite{Olah2020ZoomIn, Elhage2021MathematicalFramework, Vig2020CausalMediation, Geiger2021CausalAbstractions, geiger2025causal}. For factual knowledge in particular, prior work has argued that transformer feed-forward layers act as key-value memories and used causal tracing to localize factual associations to specific MLP computations on subject tokens, showing that these mechanisms can sometimes be edited directly \cite{geva-etal-2021-transformer, meng2022locating}. Subsequent work refined patching methodology through path-level localization, while activation-steering methods explored whether truth- or behavior-related directions can be converted into practical interventions at inference time \cite{goldowsky2023localizing, zhang2024towards, turner2023steering, li2023inference}. Our causal analysis is closest to this literature, but asks a narrower question: whether the specific correctness signal recovered by probes is merely readable, logit-causal, or strong enough to control final answer selection. 

\section{Methods}

\subsection{Task and Generation}

We evaluate short-form factual QA on SimpleQA \citep{wei2024measuring}, TriviaQA \citep{joshi2017triviaqa}, NQ-Open \citep{kamalloo2023evaluating}, PopQA \citep{mallen2023not}, and TruthfulQA \citep{Lin2022TruthfulQA:Falsehoods}. The primary model is Qwen2.5-3B, although we replicate results on Qwen2.5-7B-Instruct \citep{yang2025qwen3} and Llama-3.1-8B-Instruct \citep{grattafiori2024llama}. Our experiments use a forced-answer prompt: the model is instructed to always give an answer or its best guess. While not intended to model all real deployments, this prompt design removes abstentions and makes any confidence-correctness gap easy to observe. We also report greedy variants for several benchmarks, which expose a different failure mode: the model often abstains, leaving a smaller but more selective set of answered questions.
For each generation, we collect verbalized confidence and residual-stream states at four positions: before generation, at the first answer token, at the last answer token, and averaged over answer tokens. Appendix~\ref{app:pipeline} gives the exact prompts. Answers are graded automatically against gold aliases, with a permissive
short-answer grader that accepts exact matches, high-overlap token matches, and substring containment, and returns a binary correctness label.

\subsection{Correctness Probes}

\paragraph{Linear Probes.} For every layer and position, we fit an L2-regularized logistic regression probe on standardized residual-stream activations. Probes use a question-disjoint train/test split, so questions in the test set are not seen in
training. We primarily report AUC because the datasets are class-imbalanced: for example, SimpleQA \cite{wei2024measuring} has many more incorrect than correct answers. The probe is therefore evaluated as a ranker of correctness rather than only as a fixed-threshold
classifier. For each dataset, the headline probe is selected by held-out AUC over the grid of layers and positions. This selection can favor different positions across datasets, which is itself informative: if correctness is most readable at \texttt{prompt\_last}, the model's answer state is already present before generation; if it is most readable at \texttt{answer\_last}, correctness becomes linearly accessible only after the answer trajectory has unfolded.

\paragraph{Confidence Baselines.} We compare these linear probes against baselines computed from the generated answer: mean log-probability, minimum log-probability, first-token entropy, mean entropy, first-token margin, and mean margin. We also ask the model for a verbalized confidence score from 0 to 100 after it answers. These baselines test whether the
probe is merely rediscovering uncertainty already visible in the output
distribution. The verbalized-confidence baseline is stricter: it asks whether the model can explicitly report its own uncertainty in natural language. We also evaluate two variants that test whether the confidence result is an artifact of elicitation. In a joint prompt, the model generates both an answer and a confidence number in one response. For sampled SimpleQA generations, we compute self-consistency by majority vote over eight sampled answers. These are still external baselines: they use only what the model says, not hidden activations.

\subsection{Causal Patching}

We test causality using activation patching. For a correct run $c$ and incorrect run $w$, we replace an activation in $w$ with the corresponding activation from $c$. We then measure the change in the logit difference between correct and incorrect responses:
\[
\mathrm{LD}=\ell_{\text{gold}}-\ell_{\text{wrong}},\qquad
\Delta\mathrm{LD}=\mathrm{LD}_{\text{patched}}-\mathrm{LD}_{\text{clean}}
\]
A positive value means the patch moves the model toward the gold answer and away from the wrong answer. We patch residual streams by layer, then localize the effect to attention outputs, MLP outputs, individual attention heads, and token positions. We also run a same-question control: the correct and wrong trajectories come from different samples of the same question. In this setting the prompt is identical across trajectories, so the control patches answer-token states rather than the pre-generation prompt state. This metric is deliberately local. It does not require the entire generated answer to flip after patching; instead, it measures whether the patched state increases the model's immediate preference for the gold first token over the originally wrong first token. This makes the intervention sensitive to answer-selection mechanisms while avoiding the cost and stochasticity of full regenerated samples for every patch. We therefore use patching primarily as a localization tool: a positive rescue effect means the activation causally changes next-token evidence, not that the model's decoded answer necessarily changes.

\subsection{Directional Steering}

Finally, we test whether the probe direction can control generation. We add,
subtract, or project out multiples of the learned correctness direction during generation, including at the late MLP component that patching identifies as causally active. This tests whether a one-dimensional linear direction is a usable intervention. The comparison is important because patching inserts a full high-dimensional activation vector, whereas steering changes only one direction. If steering fails while patching succeeds, the causal information is likely distributed across a subspace or nonlinear state rather than captured by the probe normal vector alone.

\section{Experimental Setup}

\paragraph{Models.}
The main experiments use Qwen2.5-3B, a 36-layer instruction-tuned model. Replication experiments use Qwen2.5-7B-Instruct and Llama-3.1-8B-Instruct. We focus on instruction-tuned models because the target behavior is realistic question answering rather than next-token prediction on isolated statements. The 3B model is small enough for dense layer-by-layer activation collection, while the 7B and 8B models test whether our findings survive scale and different model families. All generations and patching experiments use one H100 80GB GPU with mixed precision inference.

\paragraph{Datasets.}
We generate 500-800 examples for each benchmark. For the within-question control, we sample 8 completions for each of 200 SimpleQA questions and keep questions that have at least one correct and one incorrect
completion. Probe evaluations use held-out splits over question IDs, preventing the probe from seeing the same question during training and testing. The class balance varies substantially across datasets; for instance, SimpleQA contains 83 correct and 717 incorrect answers, whereas TriviaQA contains 199 correct and 300 incorrect answers. We therefore report AUC as the primary metric and use accuracy only as a secondary summary.

\section{Results}

\subsection{Correctness Is Decodable}

Table~\ref{tab:probe_baselines} compares the best correctness probe with the best non-probe confidence baseline. Probes outperform confidence baselines on five of eight variants, with the largest gains on PopQA and greedy TriviaQA. They are not uniformly better: NQ-Open and forced TruthfulQA remain stronger under confidence baselines, suggesting that probes capture additional correctness information but do not replace output uncertainty measures.

Figure~\ref{fig:decoda_causa} shows that correctness is near chance at the embedding layer and becomes increasingly readable through the network. Peak position varies: SimpleQA peaks at \texttt{answer\_first} layer 15, while PopQA and TriviaQA rely more on late answer-token states. Thus, internal correctness is not a fixed feature; it emerges at different points depending on the task. To test whether probes merely learn question difficulty, we compare correct and incorrect samples from the same question. Among 39 SimpleQA questions with both outcomes, the best question-disjoint probe reaches AUC 0.724, while a within-question split reaches AUC 0.989. Although the latter is inflated by shared question structure, the result shows that hidden states encode answer-specific information beyond topic difficulty.

\begin{table}[t]
\centering
\caption{Best correctness probe versus the best non-probe confidence baseline.
Forced datasets use the no-abstention prompt.}
\label{tab:probe_baselines}
\begin{tabular*}{\textwidth}{@{\extracolsep{\fill}}lccc@{}}
\toprule
Dataset & Best probe AUC & Best baseline AUC & $\Delta$ \\
\midrule
NQ-Open (forced) & 0.805 & 0.833 & -0.028 \\
NQ-Open (greedy) & 0.730 & 0.817 & -0.087 \\
PopQA (forced) & 0.913 & 0.772 & +0.141 \\
SimpleQA (forced) & 0.820 & 0.785 & +0.036 \\
TriviaQA (forced) & 0.850 & 0.844 & +0.006 \\
TriviaQA (greedy) & 0.929 & 0.854 & +0.075 \\
TruthfulQA (forced) & 0.821 & 0.882 & -0.061 \\
TruthfulQA (greedy) & 1.000 & 0.942 & +0.058 \\
\bottomrule
\end{tabular*}
\end{table}

\begin{figure}[t]
    \centering
    \includegraphics[width=\textwidth]{figures/fig3_decodability_vs_causality.png}
    \caption{Correctness probe AUC by layer for SimpleQA. Correctness is generally not readable at the embedding layer, rises through the model, and often peaks in mid-to-late layers. Also shown: "rescue" effects using full activation patching from correct answers to incorrect ones. Decodability and causality do not peak at the same layer. In Qwen2.5-3B, correctness is most readable around layer 15, but patching has the largest causal effect around layer 35.}
    \label{fig:decoda_causa}
\end{figure}

\subsection{Verbalized Confidence Is a Weak Correctness Signal}

Verbalized confidence is the weakest predictor across datasets, with AUCs of 0.602 on SimpleQA, 0.785 on TriviaQA, 0.755 on NQ-Open, 0.687 on PopQA, and 0.499 on TruthfulQA. More importantly, correctness and verbal-confidence directions are nearly orthogonal in activation space. On SimpleQA, the cosine similarity between correctness and confidence directions is 0.067 for Qwen2.5-3B, 0.041 for Qwen2.5-7B, and 0.017 for Llama-3.1-8B. Incorrect-but-confident answers dominate under forced-answer prompting: 61\%, 85\%, and 72\%, respectively. Joint confidence prompting does not fix this: confidence scores correlate weakly across formats ($r=0.261$), answers match only 17\% of the time, and the incorrect-but-confident rate rises to 85\%. Likewise, self-consistency over eight SimpleQA samples gives only 8.0\% accuracy, with 78.5\% of questions never answered correctly.

\begin{figure}[h]
    \centering
    \includegraphics[width=\textwidth]{figures/fig6_confidence_dissociation.png}
    \caption{Correctness and verbal-confidence directions are nearly orthogonal. The model can internally encode that an answer is likely wrong while still producing high verbal confidence.}
    \label{fig:confidence_dissociation}
\end{figure}

\subsection{Causal Patching Localizes a Late-Layer Bottleneck}

Decodability does not imply causality, so we also attempt to patch activations from correct runs into incorrect runs. On SimpleQA with Qwen2.5-3B, residual-stream rescue is near zero early, rises after layer 22, and peaks at layer 35 with a logit-difference effect of +2.76. The best probe is at layer 15, so readability and causal influence peak at different layers. The effect is mostly logit-level rather than behavioral: full residual patching at the peak layer flips 0/60 wrong answers, MLP-only patching flips 3/60, and attention-only patching flips none. Thus, patching changes next-token evidence without reliably correcting the final answer. Figure~\ref{fig:decoda_causa} also summarizes this mismatch: probes show where correctness is readable, while patching shows where changing activations changes the computation.

\subsection{MLP Outputs Carry Most of the Causal Effect}

At the peak rescue layer, MLP-output patching dominates attention-output patching across all five Qwen2.5-3B datasets, as shown in Table~\ref{tab:patching}. Attention may still route relevant information earlier, and some individual heads show positive rescue, but the measured causal bottleneck is primarily a late-layer MLP effect rather than a sparse attention-head mechanism.

\begin{table}[t]
\centering
\small
\caption{Dataset-transfer causal patching on Qwen2.5-3B. Peak rescue layers are all
late, and MLP rescue exceeds attention rescue on every dataset.}
\label{tab:patching}
\begin{tabular*}{\textwidth}{@{\extracolsep{\fill}}lcccc@{}}
\toprule
Dataset & Peak layer & Peak rescue & MLP rescue & Attn rescue \\
\midrule
SimpleQA & L35 & 2.76 & 0.970 & 0.205 \\
TriviaQA & L34 & 7.02 & 2.620 & 0.237 \\
NQ-Open & L32 & 6.71 & 3.066 & 0.042 \\
PopQA & L34 & 7.70 & 0.725 & 0.252 \\
TruthfulQA & L35 & 11.87 & 1.190 & 0.206 \\
\bottomrule
\end{tabular*}
\end{table}

\begin{figure}[t]
    \centering
    \includegraphics[width=\textwidth]{figures/fig4_component_localization.png}
    \caption{Component localization. Late-layer MLP outputs explain substantially
    more rescue than attention outputs at the causal bottleneck.}
    \label{fig:component_localization}
\end{figure}

Same-question patching controls for topic transfer by patching between correct and incorrect samples of the same question. The largest effect again appears near layer 35, but is small (+0.467) and not statistically robust, compared with +7.676 for matched cross-question patching. We therefore treat cross-question patching as localization evidence, not a pure correctness intervention. Position-specific patching shows that \texttt{prompt\_last} carries the strongest causal signal: on SimpleQA, it peaks at layer 35 with effect +3.74. \texttt{answer\_last} also peaks at layer 35 but is smaller (+1.91), while \texttt{answer\_first} and \texttt{answer\_mean} peak near layer 30 and remain below +0.50. Much of the answer-selection state is therefore formed before the first generated token.

\subsection{Probe Directions Are Not Reliable Steering Directions}

If the probe direction were a causal control knob, steering along it should improve correctness. However, we do not observe this effect. Prompt-level steering changes SimpleQA correct rate by only +0.010 at best. Late-MLP steering at the causally active layer also fails: adding, subtracting, or projecting out the direction up to $\alpha=16$ produces no systematic improvement, with only one condition flipping 1/80
examples. This shows that a direction can be a useful readout without being a useful actuator. Full patching transfers a high-dimensional activation state, while probe steering changes only one classifier direction.

\begin{figure}[t]
    \centering
    \includegraphics[width=\textwidth]{figures/fig5_steering_negative.png}
    \caption{Directional steering is mostly ineffective. A direction can be
    predictive without being an adequate one-dimensional control knob.}
    \label{fig:steering}
\end{figure}

\begin{table}[t]
\centering
\small
\caption{SimpleQA replication across three instruction-tuned models.}
\label{tab:multimodel}
\begin{tabular*}{\textwidth}{@{\extracolsep{\fill}}lccccccc@{}}
\toprule
Model & Probe L & AUC & Rescue L & Rescue & MLP & Attn & Gap \\
\midrule
Qwen2.5-3B & L15 & 0.820 & L35 & 2.76 & 0.970 & 0.205 & +20 \\
Qwen2.5-7B & L27 & 0.850 & L27 & 5.75 & 3.068 & 0.374 & 0 \\
Llama-3.1-8B & L28 & 0.771 & L31 & 2.30 & 0.068 & 0.042 & +3 \\
\bottomrule
\end{tabular*}
\end{table}

\begin{figure}[t]
    \centering
    \includegraphics[width=0.88\linewidth]{figures/fig_multimodel.png}
    \caption{Multi-model replication. Correctness is decodable in all three models,
    and late-layer MLP effects dominate at the causal bottleneck.}
    \label{fig:multimodel}
\end{figure} 

\subsection{Replication Across Models}

The main pattern replicates on Qwen2.5-7B and Llama-3.1-8B: correctness is decodable, late-layer patches move answer evidence, and MLP outputs dominate the measured causal effect. Exact layer locations vary, so the replication is best read qualitatively rather than as evidence of identical circuits. Overall, correctness is \emph{decodable} from hidden states, partially \emph{causal} for next-token evidence, but not reliably \emph{controllable} via a single probe direction. Internal probes may therefore help monitor hallucinations, but probe directions should not be assumed to correct them.

\section{Discussion}

Our results support three claims about internal uncertainty in language models. First, correctness is not reducible to output confidence. Linear probes often outperform token-level confidence baselines, while verbalized confidence remains weak. The near-orthogonality between correctness and verbal-confidence directions suggests that stated confidence is not a faithful readout of the representation supporting answer correctness. Second, decodability does not imply causality. In Qwen2.5-3B on SimpleQA, correctness is most readable around layer 15, but causal rescue peaks around layer 35. A strong probe therefore shows where information is accessible, not necessarily where it is used to shape next-token evidence. Third, logit-level influence does not imply behavioral control. Activation patching, especially through late MLP outputs, moves next-token evidence toward the correct answer, but rarely flips the final response. Steering along the one-dimensional probe direction also fails to control correctness or verbal confidence, suggesting that the probe reads a projection of a richer causal state. These distinctions matter for hallucination monitoring. Internal probes may be useful detectors when verbal confidence is misleading, but detector directions should not be treated as intervention directions. Monitoring and control may require different representations, objectives, and levels of granularity. 

\paragraph{Limitations.} Several limitations remain. Forced-answer prompting creates many confident errors but is an artificial setting, and abstentions are high in the greedy setting (a positive sign of calibration). Following, some greedy settings have small effective test sets after abstention filtering. Additionally, probes may learn both general correctness information and dataset-specific structure, and cross-question activation patching can also transfer donor content. Thereofre, we treat it as localization evidence rather than a pure correctness
intervention. Finally, the main causal metric is a first-token logit-difference rescue score, not full answer-level behavioral correction. Future work should test more natural instruction-following settings, multi-token behavioral patching, and subspace-level interventions that preserve richer activation geometry than a single probe direction.

\paragraph{Conclusion.} We investigated whether language models have an internal signal of factual correctness and whether that signal is causal. Linear probes reveal strong correctness information in residual-stream activations, often outperforming confidence baselines and diverging sharply from verbalized confidence. Activation patching shows that late-layer MLP states can causally move next-token evidence toward the correct answer. But those shifts rarely fix the final answer, and simple steering along a learned probe direction fails to control generation or confidence. The resulting conclusion is not simply that models ``know'' when they are wrong. Rather, models expose correctness information in hidden states; some of that information lies on a causal path to answer evidence; and the causal state is richer than a one-dimensional linear probe. Internal correctness is readable, partly logit-causal, and not straightforwardly behavior-controlling.

\bibliography{references}
\bibliographystyle{plainnat}

\appendix

\section{Prompting and Pipeline Details}
\label{app:pipeline}

\subsection{Generation Prompts}

Most experiments use the forced-answer prompt below. The prompt is wrapped in the
model's chat template with a single user message and an assistant generation
prefix. This prompt is intentionally stronger than a natural deployment prompt:
it removes abstentions and creates the confidently-wrong examples needed to study
the gap between internal correctness and verbal confidence.

\begin{verbatim}
Answer the following factual question with a single short phrase. You must give your best guess; do NOT say you don't know, do NOT refuse, do NOT add commentary. Reply with the answer only.
Question: {question}
Answer:
\end{verbatim}

For the non-forced greedy variants, we use:

\begin{verbatim}
Answer the following question in one short phrase. If you do not know, say "I don't know."
Question: {question}
Answer:
\end{verbatim}

Verbal confidence is normally elicited in a separate forward pass after the model
has produced its answer:

\begin{verbatim}
Question: {question}
Your answer: {answer}
How confident are you that your answer is correct? Reply with a single integer from 0 to 100 and nothing else.
\end{verbatim}

To test whether the separate-pass design itself creates the dissociation, we also
run a joint answer-and-confidence prompt:

\begin{verbatim}
Answer the following factual question with a short phrase, then rate your confidence from 0 to 100.
You must give your best guess; do NOT say you don't know.
Format your response exactly as:
Answer: <your answer>
Confidence: <0-100>

Question: {question}
Answer:
\end{verbatim}

\subsection{Generation and Hidden-State Capture}

For the main forced-answer runs, generation is greedy with a maximum of 32 new
tokens. For the same-question and self-consistency experiments, the first sample
for each question is greedy and the remaining samples use temperature sampling
with temperature 0.7 and top-$p=1.0$.

For each generated answer, we store the prompt text, prompt length, generated
answer text, generated token IDs, answer length after dropping a trailing special
token, per-token log-probabilities, per-token entropies, and per-token margins
(top-1 minus top-2 log-probability). We then run a second forward pass over the
prompt plus generated answer with hidden-state output enabled and no KV cache.
Layer 0 is the token embedding, and subsequent indices correspond to transformer
blocks.

Residual-stream states are saved at four positions:
\begin{itemize}
    \item \texttt{prompt\_last}: the final prompt token before generation;
    \item \texttt{answer\_first}: the first generated answer token;
    \item \texttt{answer\_last}: the final generated answer token, excluding a
    trailing special token when present;
    \item \texttt{answer\_mean}: the mean hidden state over generated answer
    tokens.
\end{itemize}

These positions separate information available before generation from information
that appears during or after the answer trajectory. The hidden-state tensors are
stored as per-question files, while lightweight generation metadata is stored in
JSONL form.

\subsection{Correctness Labels}

The grader first normalizes predictions and gold answers by lowercasing,
removing punctuation, removing English articles, applying Unicode normalization,
and collapsing whitespace. An answer is marked as an abstention before grading if
it contains strings such as ``I don't know'', ``I'm not sure'', ``unknown'', or
``cannot answer''.

For non-abstaining answers, the prediction is marked correct if it satisfies any
of the following against any gold answer:
\begin{itemize}
    \item normalized exact match;
    \item the normalized gold appears in the prediction, or the prediction appears
    in the gold for short-answer cases;
    \item token F1 at least 0.6;
    \item fuzzy partial-ratio or token-set-ratio at least 85.
\end{itemize}

This grading rule is deliberately permissive because the tasks are short-form QA
with aliases and minor formatting variation. For probe and baseline evaluation,
abstentions are removed and the remaining examples are labeled correct or
incorrect.

\subsection{Probe Training and Baselines}

For each dataset, layer, and position, we standardize residual-stream activations
and train an L2-regularized logistic regression probe to predict correctness.
Unless otherwise specified, train/test splits are question-disjoint. We report
AUC as the primary metric because the datasets are class-imbalanced, and also
track accuracy, F1, log loss, and expected calibration error.

The non-probe baselines are computed directly from the generated answer: mean
log-probability, minimum log-probability, first-token entropy, mean entropy,
first-token margin, mean margin, verbalized confidence, and self-consistency where
sampled generations are available.

\subsection{Patching and Steering Implementation}

For residual-stream patching, we cache activations from a correct donor run and
replace the corresponding activation in an incorrect target run at a specified
layer and position. Component patching applies the same idea to attention outputs
or MLP outputs at the target layer. Head localization patches individual
attention-head output vectors before the output projection.

The main rescue metric is the change in first-token logit difference:
\[
\Delta =
\left[\ell(\text{gold first token}) - \ell(\text{wrong first token})\right]_{\text{patched}}
-
\left[\ell(\text{gold first token}) - \ell(\text{wrong first token})\right]_{\text{clean}}.
\]
This is a local measure of whether the intervention moves next-token evidence
toward the gold answer. It is not the same as answer-level correction, which we
measure separately by regenerating full answers after patching and grading them.

For directional steering, we add, subtract, or project out multiples of a learned
probe direction at a selected residual-stream or MLP-output site. The steering
experiments test whether a one-dimensional probe normal vector is also a usable
control direction.

\end{document}